\begin{frame}[fragile]
\frametitle{Q4-1 階乗関数の実装}
    % [c] を指定することで、左右の列の「中心軸」が垂直方向に揃います
    \begin{columns}[c]
        
        % 左側：コード
        \begin{column}{0.65\textwidth}
            % basicstyle を \tiny に下げることで、下まで途切れずに表示させます
            \lstinputlisting[
                language=C, 
                basicstyle=\ttfamily\tiny,
                numbers=left,
                numberstyle=\tiny\color{gray},
                xleftmargin=1.5em
            ]{./question4/C/1.c}
        \end{column}
        
        % 右側：アルゴリズムのポイント
        \begin{column}{0.3\textwidth}
            \small
            \textbf{アルゴリズムのポイント}
            \begin{itemize}
                \item \texttt{iszero} で終了条件を判定
                \item \texttt{pi} 関数で引数を1減らす
                \item 自己呼び出しによる再帰構造
            \end{itemize}
        \end{column}
        
    \end{columns}
\end{frame}

\begin{frame}[fragile]
\frametitle{Q4-2 数値の等価演算子の実装}
    % [c] を指定することで、左右の列の「中心軸」が垂直方向に揃います
    \begin{columns}[c]
        
        % 左側：コード
        \begin{column}{0.65\textwidth}
            % basicstyle を \tiny に下げることで、下まで途切れずに表示させます
            \lstinputlisting[
                language=C, 
                basicstyle=\ttfamily\tiny,
                numbers=left,
                numberstyle=\tiny\color{gray},
                xleftmargin=1.5em
            ]{./question4/C/2.c}
        \end{column}
        
        % 右側：アルゴリズムのポイント
        \begin{column}{0.3\textwidth}
            \small
            \textbf{アルゴリズムのポイント}
            \begin{itemize}
                \item \texttt{iszero} で終了条件を判定
                \item \texttt{pi} 関数で引数を1減らす
                \item 自己呼び出しによる再帰構造
            \end{itemize}
        \end{column}
        
    \end{columns}


\end{frame}

\begin{frame}[fragile]
\frametitle{Q4-3 数値の等価演算子の実装}
    % [c] を指定することで、左右の列の「中心軸」が垂直方向に揃います
    \begin{columns}[c]
        
        % 左側：コード
        \begin{column}{0.65\textwidth}
            % basicstyle を \tiny に下げることで、下まで途切れずに表示させます
            \lstinputlisting[
                language=C, 
                basicstyle=\ttfamily\tiny,
                numbers=left,
                numberstyle=\tiny\color{gray},
                xleftmargin=1.5em
            ]{./question4/C/3.c}
        \end{column}
        
        % 右側：アルゴリズムのポイント
        \begin{column}{0.3\textwidth}
            \small
            \textbf{アルゴリズムのポイント}
            \begin{itemize}
                \item \texttt{iszero} で終了条件を判定
                \item \texttt{pi} 関数で引数を1減らす
                \item 自己呼び出しによる再帰構造
            \end{itemize}
        \end{column}
        
    \end{columns}

    
\end{frame}

\begin{frame}[fragile]
\frametitle{Q4-4 剰余関数(mod)の実装 (1/2)}
    \begin{columns}[c]
        
        % 左側：コード前半
        \begin{column}{0.65\textwidth}
            \lstinputlisting[
                language=C, 
                basicstyle=\ttfamily\tiny,
                numbers=left,
                numberstyle=\tiny\color{gray},
                xleftmargin=1.5em,
                lastline=16 % 1行目から16行目（mod関数の終わり）までを表示
            ]{./question4/C/4.c}
        \end{column}
        
        % 右側：アルゴリズムのポイント
        \begin{column}{0.3\textwidth}
            \small
            \textbf{アルゴリズムのポイント}
            \begin{itemize}
                \item 補助関数 \texttt{lt} と \texttt{cut\_off\_sub} を活用
                \item $m < n$ になるまで再帰的に引く
                \item 引けなくなった時点の $m$ が余り
            \end{itemize}
        \end{column}
        
    \end{columns}
\end{frame}

% ---------------------------------------------------
% 2枚目：main関数の表示（17行目〜最後まで）
% ---------------------------------------------------
\begin{frame}[fragile]
\frametitle{Q4-4 剰余関数(mod)の実装 (2/2)}
    \begin{columns}[c]
        
        % 左側：コード後半
        \begin{column}{0.65\textwidth}
            \lstinputlisting[
                language=C, 
                basicstyle=\ttfamily\tiny,
                numbers=left,
                numberstyle=\tiny\color{gray},
                xleftmargin=1.5em,
                firstline=17 % 17行目（main関数）から最後までを表示
            ]{./question4/C/4.c}
        \end{column}
        
        % 右側：後半のポイント
        \begin{column}{0.3\textwidth}
            \small
            \textbf{メイン関数の処理}
            \begin{itemize}
                \item 標準入力からの値の受け取り
                \item 0除算エラーのハンドリング（\texttt{-1} が返った場合）
                \item 結果のフォーマット出力
            \end{itemize}
        \end{column}
        
    \end{columns}
\end{frame}

\begin{frame}[fragile]
\frametitle{Q4-5 数値の等価演算子の実装}
    % [c] を指定することで、左右の列の「中心軸」が垂直方向に揃います
    \begin{columns}[c]
        
        % 左側：コード
        \begin{column}{0.65\textwidth}
            % basicstyle を \tiny に下げることで、下まで途切れずに表示させます
            \lstinputlisting[
                language=C, 
                basicstyle=\ttfamily\tiny,
                numbers=left,
                numberstyle=\tiny\color{gray},
                xleftmargin=1.5em
            ]{./question4/C/5.c}
        \end{column}
        
        % 右側：アルゴリズムのポイント
        \begin{column}{0.3\textwidth}
            \small
            \textbf{アルゴリズムのポイント}
            \begin{itemize}
                \item \texttt{iszero} で終了条件を判定
                \item \texttt{pi} 関数で引数を1減らす
                \item 自己呼び出しによる再帰構造
            \end{itemize}
        \end{column}
        
    \end{columns}

    
\end{frame}